%───────────────────────────────────────────────────────────────────────────────
% EQUATION-STYLES.TEX — Coloreado semántico de ecuaciones (taxonomía híbrida)
%───────────────────────────────────────────────────────────────────────────────
%
% Sistema de "syntax highlighting" para matemáticas que combina tres ejes:
%
%   1. ROL SINTÁCTICO (estable entre contextos)
%      Operador, función, diferencial, dominio — siempre tienen el mismo color.
%
%   2. ESTATUS EPISTÉMICO (la dimensión que la tipografía NO codifica)
%      Incógnita vs dato/parámetro fijo vs constante universal.
%      En f(x;θ): en física x es incógnita y θ dato; en ML es al revés.
%      El color hace explícita esta inversión.
%
%   3. TIPO MATEMÁTICO (refuerzo sutil de la tipografía)
%      Vector, tensor/matriz — la tipografía (bold, doble barra) ya lo indica,
%      el color refuerza cuando el contexto lo amerita.
%
% Regla de uso: si el símbolo tiene un rol sintáctico claro (operador, función,
% diferencial, dominio), usa ese macro. Si no, pregúntate: ¿es conocido o
% desconocido? Usa \eqk o \equ. Si además quieres reforzar que es vector o
% tensor, usa \eqvec o \eqten. Si ningún color aporta información, déjalo
% sin colorear — el sistema es opt-in, no obligatorio.
%
% Toggle global: \eqcolorstrue / \eqcolorsfalse
%───────────────────────────────────────────────────────────────────────────────

% ─── Toggle global ───────────────────────────────────────────────────────────
\newif\ifeqcolors
\eqcolorstrue                         % \eqcolorsfalse para versión B/N

% ═════════════════════════════════════════════════════════════════════════════
% MAPA SEMÁNTICO → COLOR
% ═════════════════════════════════════════════════════════════════════════════
%
% ── Eje 1: Rol sintáctico (estable) ──────────────────────────────────────────
%
%   Macro      Color       Cubre
%   ─────      ─────       ─────
%   \eqop      Blue        Operadores: ∫, ∑, ∏, ∂, ∇, ∇², ∇×, lím, Ĥ
%   \eqfn      Green       Funciones/mapas: f(·), L(θ), σ(·), V(r), ReLU
%   \eqdm      Teal        Diferenciales/medidas: dx, dt, dW, dS, dΩ
%   \eqdom     Sapphire    Dominios/conjuntos: Ω, ∂Ω, ℝⁿ, H, S²
%
\colorlet{eqop}{CtpBlue}
\colorlet{eqfn}{CtpGreen}
\colorlet{eqdm}{CtpTeal}
\colorlet{eqdom}{CtpSapphire}

% ── Eje 2: Estatus epistémico ────────────────────────────────────────────────
%
%   Macro      Color       Cubre
%   ─────      ─────       ─────
%   \equ       Peach       Incógnita / buscado / aprendido
%                          Lo que se quiere encontrar, optimizar, resolver.
%                          Física: estado x, campo ψ, distribución p(x).
%                          ML: pesos θ, predicción ŷ, política π(a|s).
%
%   \eqk       Mauve       Conocido / dato / fijado
%                          Lo que se asume dado o se mide.
%                          Física: parámetros θ, condiciones de frontera.
%                          ML: datos x, labels y, hiperparámetros η.
%
%   \eqcon     Lavender    Constante universal / estructural
%                          Invariante del contexto: ℏ, k_B, π, c, ε₀, i.
%                          También factores numéricos fijos (1/2, 2π).
%
\colorlet{equ}{CtpPeach}
\colorlet{eqk}{CtpMauve}
\colorlet{eqcon}{CtpLavender}

% ── Eje 3: Tipo matemático (refuerzo sutil) ─────────────────────────────────
%
%   Macro      Color       Cubre
%   ─────      ─────       ─────
%   \eqvec     Red         Campos vectoriales: u, E, B, J, f
%                          (la tipografía bold ya lo indica; el color refuerza)
%
%   \eqten     Rosewater   Tensores / matrices: σ_ij, g_μν, W, Σ
%                          (para cuando distinguir rango es importante)
%
\colorlet{eqvec}{CtpRed}
\colorlet{eqten}{CtpRosewater}

% ═════════════════════════════════════════════════════════════════════════════
% MACROS
% ═════════════════════════════════════════════════════════════════════════════
\makeatletter

% ─── Comando interno ─────────────────────────────────────────────────────────
\newcommand{\@eqc}[2]{%
  \ifeqcolors\textcolor{#1}{#2}\else#2\fi
}

% ── Eje 1: Rol sintáctico ───────────────────────────────────────────────────
\newcommand{\eqop}[1]{\@eqc{eqop}{#1}}       % operador
\newcommand{\eqfn}[1]{\@eqc{eqfn}{#1}}       % función / mapa
\newcommand{\eqdm}[1]{\@eqc{eqdm}{#1}}        % diferencial / medida
\newcommand{\eqdom}[1]{\@eqc{eqdom}{#1}}      % dominio / conjunto

% ── Eje 2: Estatus epistémico ───────────────────────────────────────────────
\newcommand{\equ}[1]{\@eqc{equ}{#1}}          % incógnita / buscado
\newcommand{\eqk}[1]{\@eqc{eqk}{#1}}          % conocido / dato
\newcommand{\eqcon}[1]{\@eqc{eqcon}{#1}}      % constante universal

% ── Eje 3: Tipo matemático (refuerzo) ───────────────────────────────────────
\newcommand{\eqvec}[1]{\@eqc{eqvec}{#1}}      % vector
\newcommand{\eqten}[1]{\@eqc{eqten}{#1}}      % tensor / matriz

% ═════════════════════════════════════════════════════════════════════════════
% MACROS COMPUESTOS
% ═════════════════════════════════════════════════════════════════════════════

% Derivada parcial coloreada
% Uso: \eqpdv{ρ}{t}  →  ∂ρ/∂t con operador + incógnita/variable
% El primer argumento se colorea como incógnita, el segundo como conocido
% (la variable independiente es la coordenada, un dato del problema).
\newcommand{\eqpdv}[2]{%
  \eqop{\frac{\partial \equ{#1}}{\partial \eqk{#2}}}%
}

% Gradiente coloreado
% Uso: \eqgrad{u}  →  ∇u
\newcommand{\eqgrad}[1]{%
  \eqop{\nabla}\equ{#1}%
}

% Divergencia coloreada
% Uso: \eqdvg{\mathbf{u}}  →  ∇·u
% (nombre \eqdvg para no colisionar con \eqdiv si se usa \div de physics)
\newcommand{\eqdvg}[1]{%
  \eqop{\nabla \cdot} \equ{#1}%
}

% SDE coloreada
% Uso: \eqSDE{X}{-\theta X}{\sigma}
% Estado = incógnita, drift = función, coeficiente difusión = conocido
\newcommand{\eqSDE}[3]{%
  \eqdm{d}\equ{#1} = \eqfn{#2}\,\eqdm{dt} + \eqk{#3}\,\eqdm{dW}%
}

\makeatother